\documentclass[12pt]{article}
\usepackage[utf8]{inputenc}
\usepackage[T1]{fontenc}
\usepackage{hyperref}
\title{Khua Language Resources for K-12 Curriculum Development}
\author{Compiled by Claude Sonnet 5.5 (claude-sonnet-5-5) for Andr\'as Kornai}
\date{}
\begin{document}
\maketitle

\section{Introduction}
Khua (ISO 639-3 xhv; Glottocode khua1243) is an Austroasiatic (Mon-Khmer) language of Vietnam and Laos. Search results place it in the West Katuic group, closest to Bru, So and Mangkong. Joshua Project reports about 4,900 Khua in Laos (northern Bolikhamxay/Khammouan area, Bouarapha District) and about 1,300 in Vietnam, where they are counted within the Bru-Van Kieu minority. The population figures are Joshua Project's and may differ from other sources. Khua should not be confused with the Cua (a Bahnaric-speaking group) or with Katua (kta), a separate language.

\section{Primary Educational Resources}
I found no primers, textbooks, school programmes or orthography for Khua in my searches. The reference description says it is not known to be taught in schools, and I found nothing that contradicts this.

\section{Linguistic and Reference Materials}
\begin{itemize}
\item Joshua Project, language page for Khua (xhv): \url{https://joshuaproject.net/languages/xhv}
\item Joshua Project, Khua in Laos: \url{https://joshuaproject.net/index.php/people_groups/12668/LA}
\item Joshua Project, Khua in Vietnam: \url{https://joshuaproject.net/index.php/people_groups/12668/VM}
\item Ethnologue subgroup page (Bru subgroup context): \url{https://ethnologue.com/subgroup/1755/}
\item Kiwix copy of Wikipedia, ``Khmuic languages'' (neighbouring Austroasiatic group, not Khua itself): \url{https://kiwix.lema.org/wikipedia_en_all_mini_2024-02/A/Khmuic}
\end{itemize}
No dictionary, grammar or corpus for Khua turned up. The Joshua Project pages are demographic and missionary profiles, not linguistic descriptions.

\section{Religious and Community Resources}
I found no scripture translations or community organisations specific to Khua. The Joshua Project pages are the only people-group profiles I saw.

\section{Digital and Audio Resources}
I found no audio, video or software resources for Khua.

\section{Curriculum Development Context}
Documentation is thin, so a Khua curriculum would have to start with community-led documentation. That means recording speakers, agreeing an orthography and building a basic word list. Related Katuic languages such as Bru, with better documentation, could offer models, but I did not examine them in this search.

\end{document}